\documentclass[11pt]{article}
\usepackage[T1]{fontenc}
\usepackage[utf8]{inputenc}
\usepackage{lmodern}
\usepackage[margin=1in]{geometry}
\usepackage{amsmath,amssymb,amsthm,mathtools}
\usepackage{enumitem,booktabs,longtable,array}
\usepackage{microtype}
\usepackage[hidelinks]{hyperref}
\newtheorem{principle}{Candidate principle}[section]
\newtheorem{target}{Research target}[section]
\theoremstyle{definition}
\newtheorem{definition}{Working definition}[section]
\newtheorem{control}{Baseline control}[section]
\newcommand{\Ob}[2]{\subsection*{#1\quad #2}\addcontentsline{toc}{subsection}{#1\quad #2}}
\newcommand{\Ans}{\mathsf{Ans}}
\newcommand{\Real}{\mathsf{Real}}
\newcommand{\Obs}{\mathsf{Obs}}
\newcommand{\Cost}{\mathsf{Cost}}
\newcommand{\Ctx}{\mathsf{Ctx}}
\newcommand{\Law}{\mathsf{Law}}
\newcommand{\Poss}{\mathsf{Poss}}
\newcommand{\Cat}{\mathbf{Cat}}
\newcommand{\Set}{\mathbf{Set}}
\newcommand{\RR}{\mathbb R}
\newcommand{\NN}{\mathbb N}
\newcommand{\supp}{\operatorname{supp}}
\newcommand{\id}{\operatorname{id}}
\newcommand{\one}{\mathbf 1}
\setlist{leftmargin=*,itemsep=3pt,topsep=5pt}
\setlength{\parskip}{4pt}
\setlength{\emergencystretch}{2em}
\title{Categorical Phenomenology Programme Contract\\[4pt]\large Formation, warrant, and transformation of structured understanding}
\author{Proposal for adoption}
\date{4 November 2026}
\begin{document}
\maketitle

\begin{abstract}
The intellectual prize is a mathematics of how a field of encounters becomes a field of intelligible objects, operations, and laws, and how that organization can be justified, used, and transformed without concealing the conditions that made it possible. The programme must reach beyond recovering states inside a supplied model. It will investigate the formation of scientific and mathematical theories, the interplay of different kinds of warrant, and the composition of heterogeneous accounts that need not be descriptions of one recoverable hidden object.

This contract commits the programme to positive constructions, substantive comparisons, an integrated account of changing languages and practices, and sharp limits. It makes access and use, constitutive heterogeneity, and immanent transformation explicit research hypotheses rather than compulsory definitions. No categorical technology is selected in advance. The desired outcome is a body of mathematics that can determine what is constructed, what is justified, what can be done with it, what must change when those uses fail, and which of these determinations can be transported between genuinely different settings.

The obligations below are proposed commitments, not achieved results. This is a standalone replacement proposal for consideration, not an adopted replacement or authorization to execute the whole programme. Existing results retain their scope and credit. The carry-forward and change ledger makes the proposed amendments to the earlier contract explicit; no positive obligation is silently discharged by a counterexample, a new vocabulary, or this document itself.
\end{abstract}
\tableofcontents

\section{The intellectual undertaking}
\label{sec:undertaking}

A theory does more than compress records. It determines what can count as an object, a difference, an explanation, an experiment, or a proof. It can introduce a quantity that makes a previously inarticulate regularity expressible; an operation that makes new evidence obtainable; or an abstraction that reveals that several old constructions were instances of one law. It can also conceal distinctions that become decisive when the theory is used. Categorical phenomenology will study these changes in organized intelligibility as mathematical problems with explicit inputs, conclusions, and warrants.

The programme's ambition is to make the formation of such structure tractable without pretending that structure arises from nothing. Instruments, inferential practices, languages, objectives, and mathematical commitments participate in what can be formed. Their participation is not merely a nuisance to be listed in a preamble. Where it affects the answer, it becomes part of the explanatory problem. We will determine how varying those conditions changes the resulting concepts, laws, realizations, and possibilities of action, and identify circumstances in which a principle survives that variation.

The resulting theory must be capable of discovering a common mechanism and of discovering that a common mechanism is the wrong demand. A proof, a likelihood, an intervention, a universal construction, and a practical decision do not become interchangeable because they fit into a large category. Equally, their differences need not prevent rigorous cooperation. An important achievement would be a law explaining how a deductively certified construction, statistically supported identification, and operationally licensed intervention can jointly sustain a conclusion that none supports alone.

The programme will therefore pursue a connected positive achievement: \emph{construct and justify new structure; explain relations among unlike constructions; account for the formation and revision of languages and theories; and prove when the resulting accounts can be used together}. Correctly identifying the limits of these achievements is part of the mathematics. A limit does not substitute for all of them.

This is not a promise to automate arbitrary discovery, derive relevance from uninterpreted experience, or select a final ontology. Nor is it a programme to build ever more elaborate models of one convenient example. A local model earns its role by deciding a conceptual question. The central concepts must remain capable of being changed by what those investigations reveal.

\paragraph{Status of the name.}
The mathematical subject itself is under investigation. Connections to philosophical phenomenology may motivate or clarify a particular claim, but a historical genealogy is not an acceptance gate. First-person experience, intentionality, and embodiment are not claimed as mathematical consequences of observational reconstruction. Any result explicitly about them must supply its own interpretation and warrant.

\section{The subject and its competing hypotheses}
\label{sec:subject}

The working subject is the formation, justification, organization, and transformation of structure under conditions of encounter, interpretation, and use. This is an orientation from which precise problem classes will be constructed, not a universal definition established by fiat. The programme must test whether it identifies a substantive mathematical family rather than redescribing all mathematics.

\begin{principle}[Dependence on access and use]
For a substantial class of formation problems, the distinctions and structures that can be warranted depend jointly on how the subject matter can be encountered and what conclusions or operations the resulting account must support. An adequate comparison must preserve enough of both to transport the intended warrant.
\end{principle}

Here access may be experimental, operational, inferential, or constructive, with different meanings in different problems. The claim is not that a proof is a measurement or that every mathematical construction is secretly a probe. Use can include explanatory organization, proof construction, intervention, prediction, or coordination; it is not restricted to external instrumental benefit. This candidate must be compared with formulations organized by consequence and interpretation, two-sided admissibility, dynamical continuation, or local compatibility. Its value must come from distinctions and theorems it makes possible, not from the flexibility of its wording.

\begin{principle}[Constitutive heterogeneity]
Some phenomenological problems are constituted by interactions among different forms of presentation, evidence, operation, and justification. Their solution cannot always be obtained by translating all components into one homogeneous semantics and then forgetting how the translation was made.
\end{principle}

This hypothesis requires more than a heterogeneous application. The programme must exhibit a class in which erasing a specified difference loses an intended answer, an admissible operation, a warrant, or a compositional guarantee, and construct a smaller adequate account that retains precisely the relevant difference. It must also establish cases in which heterogeneity \emph{can} be removed. Preserving every detail is no more explanatory than erasing every difference.

\begin{principle}[Formation can be immanent and partly unrecognized]
Practices, interactions, and transformations can alter the distinctions or operations available within a situation before the change has been explicitly represented or recognized. A later reflective reconstruction can explain that development without having caused or accompanied every stage of it.
\end{principle}

The programme will distinguish transformations of a subject matter or practice, transformations of its representations, and the construction of tools for understanding either. Recognition, explicit diagnosis, and certified revision are important possible events, not prerequisites for all change. No universal pipeline of observation, recognition, revision, and validation is imposed. ``Dialectical'' will denote an explicitly studied relation among internal tensions, practical transformation, and reorganization when one can be given; it will not serve as an unexplained transition operator.

A fixed ambient dynamical or state-space formalism may faithfully represent changes of vocabulary, goals, or capabilities. Its being a model does not disqualify it. The issue is whether its construction explains the claimed formation and dependence rather than merely storing every possible future organization in a supplied carrier.

\begin{principle}[Integration has mathematical content]
When structural, deductive, epistemic, interventional, prescriptive, and computational conclusions are used together, their compatibility is an additional mathematical question. Locally valid conclusions need not compose into a valid explanation or action.
\end{principle}

\begin{principle}[Unity must earn its scope]
A common construction, an interoperating family of doctrines, and an obstruction to stronger unity are competing possible outcomes. Their scope is determined by positive constructions and exact comparisons, not by advance preference for a universal formalism or irreducible plurality.
\end{principle}

\Ob{P1}{Determine a substantive mathematical subject}
Develop precise problem classes realizing these candidates, and prove that at least one proposed boundary has mathematical consequences: an invariant, representation criterion, obstruction, or classification that distinguishes its class from an unrestricted ambient category. Supply genuine inclusions and exclusions. In particular, distinguish a problem about the dependence and formation of warranted structure from an arbitrary construction relabelled as phenomenology, and explain why mathematical theory formation is not excluded merely because it is not empirical sampling.

Compare at least the access/use, logical-interpretive, and contextual-realization candidates on overlapping questions. A successful foundational result must change the organization of the programme: determine an essential dependence, a principled division, or a reusable construction. A slogan, taxonomy, or permissive tuple does not meet P1. No result is required to prove that every future legitimate example has already been captured.

\Ob{P2}{Explain which questions make a representation adequate}
Develop a theory of question-relative adequacy that permits objects, identified sets, proofs, constructions, predictions, and admissible actions as genuinely different answers. Determine how changes of question and intended use alter the adequate representation, and which changes leave a common sufficient fragment. Give a positive construction of such a fragment on a substantial class, together with an obstruction where no single best fragment exists.

Relevance is not to remain an unexplained parameter everywhere. Prove at least one result deriving or constraining a relevant family of questions from independently specified tasks, inferential commitments, interaction structure, or invariance requirements. The supplied objective or commitment remains visible. A prescribed objective does not become an objective forced by evidence.

\section{Answers, warrants, and the conditions of adequacy}
\label{sec:warrant}

The following small interface fixes distinctions needed by several obligations. It is not a proposed universal ontology. A logical formulation may implement it through judgments and interpretations; an experimental formulation through protocols and uncertainty; a geometric one through local data and descent.

\begin{definition}[Answer adequacy on a fixed problem class]
Let $X$ be a class of source situations. For each question $q\in Q$, independently specify a required domain $D_q\subseteq X$ and admissible-answer sets $\Ans_q(x)\subseteq Y_q$ for $x\in D_q$. A representation $r:X\to R$ is answer-adequate for $Q$ if each question has an answer procedure or semantic selector $a_q$ defined on every point of $r(D_q)$ such that
\[
 a_q(r(x))\in\Ans_q(x)
\]
for every $x\in D_q$. A theorem must distinguish semantic existence of $a_q$, an effective construction of it, and the warrant for accepting its output. Partial adequacy must declare its required domain independently rather than removing difficult cases after the construction. If $\Ans_q(x)=\varnothing$ on a required source case, no adequate representation exists for that requirement. An obstruction or abstention is an answer only where the original question permits it. If deciding applicability is itself required, it is a further preserved question.
\end{definition}

For exact uniquely answered questions this requires equality of answers on each relevant fibre of $r$. For multivalued questions it requires a common admissible answer, not equality of entire answer sets. In a finite set-based control, answer selection at $z\in r(D_q)$ is possible precisely when
\[
 \bigcap_{x\in r^{-1}(z)\cap D_q}\Ans_q(x)\ne\varnothing.
\]
For arbitrary families, simultaneous selection carries its usual choice or constructivity qualifications. This elementary criterion is a baseline, not the programme's new sufficiency theorem. It says nothing by itself about obtaining $r(x)$, computing an answer, or supplying a valid certificate. Asking only questions manufactured to be constant on a chosen quotient would make adequacy circular.

An identified set or confidence region can itself be one value in $Y_q$, with its own inclusion, coverage, and sharpness requirements. This differs from $\Ans_q(x)$, the set of all acceptable answers. The common-answer criterion does not require selecting one true model when the task asks for a valid region. Jointly required answers may also obey compatibility conditions not captured by independent questionwise selection. A decoder or warrant transformer may not secretly consult the inaccessible source state that its representation is supposed to replace.

The common $Y_q$ is justified only on a class with a shared answer interpretation. Heterogeneous answer types require explicit, possibly partial decoders, translations, or correspondences and a proof of the warrant they preserve. A disjoint union of types is not by itself a meaningful comparison.

There need not be a coarsest adequate quotient. For three situations with answer sets $\{a\}$, $\{a,b\}$, and $\{b\}$, both partitions $\{1,2\}\mid\{3\}$ and $\{1\}\mid\{2,3\}$ admit sound choices, but their common coarsening does not. This elementary control preserves the maximal-congruence target where its stronger exact-answer hypotheses hold, while ruling out its promotion to a universal adequacy principle.

\paragraph{Warrant is typed.}
A deductive result has premises, a consequence relation, and a proof or semantic argument. An exact reconstruction has a declared ideal law and construction assumptions. A statistical result has an evidence-generating model and a procedure-level coverage, risk, or convergence statement. An interventional result has operational or causal assumptions relating manipulation to outcomes. A prescription has an objective, constraints, and an optimality, regret, admissibility, or dominance claim. A computational result has an input representation, an algorithm or reduction, and a cost model.

A confidence statement is not a proof that one realized assertion is true. Its mathematical warrant may be a proof that a procedure covers a parameter with probability at least $1-\delta$ under stated assumptions. A translation must preserve or explicitly change that guarantee, rather than treating statistical confidence as another truth value. Likewise a correct proof in a formal language does not validate the language's empirical interpretation.

\Ob{P3}{Construct a calculus of heterogeneous warrant}
Construct a calculus of claims and their justification in which at least deductive, exact structural, and statistical or operational warrants can be combined without conflation. The calculus must derive a nontrivial composite claim, identify the extra compatibility premises, and support transport or revision of the appropriate certificates. It must distinguish sound transport, reflection, conservative extension, loss of precision, and loss of applicability.

The main positive theorem must say when an assembled justification is valid in the source problem, rather than merely internally derivable in a target notation. Include a counterexample in which individually valid components produce an invalid joint conclusion, and a repaired positive class. Possible failures include a statistical guarantee invalidated by adaptive selection, a formally valid abstraction used outside its interpretation, and locally consistent claims whose joint use silently adds a dependence assumption. A generic product of certificate types does not meet this obligation.

\Ob{P4}{Relate semantic possibility to obtainable evidence}
Construct an exact and a quantitative account of which question-relevant distinctions can be acquired through declared operations. Distinguish available records, currently executable protocols, complete ideal laws on those protocols, and hypothetical larger law domains. Address sequential or destructive access, rather than assuming every separately available experiment can be jointly performed.

Prove an acquisition or identification theorem with explicit progress and termination or convergence conditions on a nontrivial class; pair it with an indistinguishability or access obstruction. At least one result must show how the structure of the desired answer changes the required experiment, rather than reconstructing every hidden datum first. At least one must address the gap between semantic sufficiency and executable sufficiency. Existing finite belief-state and replay results are dependencies, not sufficient completion by repetition. Statistical claims require their own sampling, selection, confidence, and stability assumptions.

\Ob{P5}{Construct structure rather than accommodate it}
Prove a new positive construction theorem deriving objects, concepts, relations, operations, invariants, realizations, or an identified family from independently stated presentations and commitments. State precisely what is generated rather than supplied; prove its universal, extremal, inferential, operational, or quantitative property; and show how admissible input changes alter the result.

The six inherited nonvacuity conditions remain binding: specified output and use; visible supplied structure; independent input hypotheses; a positive nontrivial family; counterfactual dependence and relevant invariance; and explicit warrant, ambiguity, and strongest-source comparison. An independent abstract hypothesis need not be a finite algorithm, but a claimed finite or effective criterion must actually be one. A chosen useful structure is permitted, provided its selection criterion and alternatives are exposed. Neither a preinstalled latent architecture nor an image factorization with a new name satisfies the new-result commitment.

\Ob{P6}{Explain comparison across unlike mechanisms}
Prove a positive bridge across substantively different construction mechanisms, with source-level meaning, exact direction, a derived consequence, and a sharp boundary. Preserve the earlier incidence/closure, temporal/predictive, and native-factorization comparison as a named obligation, while extending the programme beyond that trio. A comparison may be an equivalence, directed interpretation, adjunction, span, simulation, or quantified loss, but it must preserve the particular answers and warrants claimed.

The common result must do work that the separate constructions do not already supply: derive a new output, transport a nontrivial theorem or experiment, expose a necessary distinction, or determine a previously unresolved compatibility condition. A new theorem-level claim requires a residual beyond the strongest relevant predecessors. An explanatory synthesis can be credited independently without claiming priority. A failed inverse is not the absence of every directed comparison, and an impossibility theorem must quantify over its allowed translation class.

\section{Formation of languages and theories}
\label{sec:formation}

Formation is a central positive commitment, not an optional application of a completed reconstruction theory. The programme will distinguish construction inside a language, identification within a hypothesis family, selection under an objective, introduction of new concepts or operations, and revision of the language or its interpretation. Several can occur together, but a theorem about one does not establish the others.

The crucial advance is to make some conditions of theory formation mathematical without installing the sought theory in the input. Inputs may include grammars, admissible construction rules, proof systems, instrument semantics, and inductive commitments. The target is not assumption-free invention. It is a theorem explaining why a meaningful new organization arises, what consequences it enables, which alternative organizations remain, and how its legitimacy depends on those inputs.

\Ob{P7}{Produce a scientific theory formation theorem}
For a nontrivial class of scientific inference problems, construct or identify a change in explanatory structure that cannot be achieved by parameter estimation within the original fixed language. The change must introduce or reorganize a law, variable, mechanism, operation, or measurement interpretation. Prove a consequence that was not simply used as a defining constraint, and state what evidence or operation could test it.

The theorem must relate the formation step to independently specified empirical or operational constraints and distinguish rival explanations such as a new mechanism, altered measurement, contamination, or changed background conditions where those alternatives are admitted. It must retain an identified family when the evidence does not select one explanation. A positive result may identify an equivalence class or a bounded language extension rather than a unique ontology.

Success requires a mathematically justified passage from an inadequate or incomplete organization to a useful richer or differently organized one, a justified account of what survives, and an ambiguity or failure boundary. A supplied rational recurrence replacing a supplied linear model is a valuable control but not the whole obligation. The next theorem must address a family-level formation rule, an input-level selection criterion, or another additional mechanism beyond that example. Synthetic data or ideal-law input must be labelled as such; predicted consequences are not reported as experimentally confirmed.

\Ob{P8}{Produce a mathematical theory formation theorem}
Develop a formation principle in which mathematical examples, constructions, proofs, counterexamples, or compatibility constraints warrant new definitions, sorts, operations, axioms, or universal properties. Construct the resulting theory or class of theories and prove an explanatory consequence: a common proof, a new transfer, a classification, a previously unavailable construction, or a precise independence result.

The input may specify an admissible language of theory transformations, but may not designate the final theory and call its restatement formation. Show how the mathematical material constrains the introduced organization and why a materially different choice fails a declared explanatory requirement. Distinguish a definitional extension, a conservative but expressive extension, a nonconservative axiom addition, a Morita-style change of presentation, and a change of consequence relation when relevant. New axioms require a separate justification; they do not become theorems of the old theory by being useful.

At least one theorem must show that the proposed organization makes a proof or construction available across distinct mathematical settings while preserving their relevant interpretations. The source of its warrant is deductive or constructive, not a fictional sampling process. Compare against categorical logic, presentation theory, classifying constructions, and existing heterogeneous logic and theory-formation methods before asserting novelty.

\Ob{P9}{Construct language change with controlled interpretation}
Give a theory of transformations that can change signatures, questions, operations, or inference rules, rather than only extending a table in a fixed language. Construct the interpretation or comparison between the old and new accounts, identify its domain, and prove which claims, proofs, empirical commitments, and operational capabilities survive.

A principal target is a criterion separating: extra capacity in the same language; a coordinate change; a definitional extension; an expressive change; a change of empirical interpretation; and a change of task. At least two of these must interact nontrivially in the positive theorem. Prove a comparison for successive changes and identify when their order matters. A language change may preserve one warrant and invalidate another; the theory must locate that difference rather than demand universal conservativity.

\Ob{P10}{Explain formation before and after recognition}
Construct a mathematical account in which an evolving practice or interaction can change effective distinctions, operations, or inferential possibilities before a reflective description of that change is available. Prove a positive relation between the underlying transformation and a later reconstructed account, with explicit loss or partiality. Also exhibit a class in which some change remains unrecognized within the available reflective resources.

This obligation must not be satisfied by adding a Boolean ``recognized'' flag to a fixed model. The transformation must alter a substantive space of possible questions, operations, or justified conclusions; the reflective account must recover or explain something about that alteration. Investigate at least one genuinely immanent source of revision, such as incompatible internal commitments, failure of compositional use, or a construction that makes its previous language inadequate. No theorem is required to make every transformation conscious, explicit, diagnosable in finite time, or representable by the same revision operator.

\Ob{P11}{Explain stabilization without freezing the theory}
Prove conditions under which a developing organization achieves a specified form of stability while remaining revisable. Distinguish persistence of an invariant, interpretability under language change, convergence of a learning or inquiry process, robustness to perturbation, and persistence of explanatory or operational adequacy. Establish a positive stability result and an instability or bifurcation boundary that determines what kind of revision is needed.

At least one result must permit the carrier, vocabulary, or presentation to change while a substantive organized relation persists. Conversely, exhibit a stable fixed representation that ceases to be adequate when the question or context changes. Stabilization must not be inferred from repeated naming, parameter convergence alone, or preservation of every old formal object. Any optimality criterion for selecting a stable theory remains a supplied or separately justified commitment.

\Ob{P12}{Make the programme reflexively informative}
Use the developing mathematics to explain a substantive change in the programme's own tools or problem formulations. A result must show why an earlier organizing account lost a relevant distinction, how a revised formulation recovers it, and what additional theorem or construction becomes available. The account must distinguish the development itself from its later reconstruction.

This is a bounded mathematical self-application, not a demand for a final self-justifying metatheory. It must not claim that its own axioms are established merely by its success or that every practical transformation has been formally represented. A successful instance could explain why a law-only account requires possibility data, why a static testing account requires executable access, or why an observation-preserving language extension does not preserve a geometric witness structure. Repeating those known examples is preparatory; the new obligation is a further, proved account of formulation change and its consequence.

\section{Heterogeneous composition and evolving inquiry}
\label{sec:composition}

The programme must explain interactions among its dimensions, not assign each dimension a separate chapter and call the collection integrated. Structural formation can change the available proof language. A proof can justify an experiment's interpretation. The experiment can change the system as well as its description. A prescribed policy can determine which evidence becomes available, and the resulting evidence can invalidate the abstraction used to prescribe it. The mathematics must expose these dependencies at the level at which a conclusion is claimed.

\Ob{P13}{Construct global use from partial accounts}
Develop constructive local-to-global results for heterogeneous descriptions, constraints, theories, or behaviors. Distinguish overlap compatibility, global existence, uniqueness, effective gluing, and adequacy for the intended use. Prove a positive assembly theorem without supplying the desired global witness in advance, an explicit obstruction, and a useful weaker construction when a global account is unavailable.

In particular, determine when a family of partial accounts or directed translations supports a task even though no single global model exists in the selected doctrine. Conversely, identify a question whose answer genuinely requires additional joint information. Correlations, time alignment, scale changes, and semantic identifications must be supplied, inferred with warrant, or kept unresolved. Sheafification or free completion cannot be passed off as a source-realizable world. Existing descent and contextuality results are controls; the new theorem must connect assembly to formation, warrant, revision, or operational use in a way they do not already settle.

\Ob{P14}{Derive contextual adequacy and replacement}
Construct input-level criteria under which a representation or theoretical account adequate in isolation remains adequate under its declared sequential, parallel, adaptive, or interventional uses. Derive the closure or congruence rather than assume the desired replacement theorem. State the exact class of clients, substitutions, or contexts, the semantics they inspect, and what the conclusion preserves.

The positive result must construct a missing distinction, an adequate interface, a justified closure, or a controlled repair where present-response equality is insufficient. It must retain decomposition witnesses, occurrence structure, symmetries, or noninvertible comparison cells whenever those affect a source question. If feedback or trace is claimed, prove definedness and the required preservation law. An unbounded context class can have a finite certificate; a finite list of tested contexts is not automatically such a certificate.

\Ob{P15}{Establish laws for repeated and interacting revisions}
Construct comparisons for successive changes of evidence, access, interpretation, language, doctrine, objective, and implementation. These are not all inclusions or Bayesian updates. Prove conditions for coherent claim reuse and for retraction, weakening, or replacement where reuse fails.

A central result must decide an actual interchange problem: forming structure before a change versus forming it after the change; composing locally revised accounts versus revising their assembly; or introducing a language before an intervention versus interpreting its aftermath. Produce the correctly typed equality, invertible comparison, directed correction, span, error bound, or obstruction. Exhibit an order effect that changes a substantive answer and a positive class in which an explicit retained distinction restores a useful law. Do not enforce commutation where the order encodes a real causal or historical difference.

\Ob{P16}{Generate new means of inquiry}
Construct a formation process whose output makes a new question, proof method, experiment, instrument specification, or operation available. Prove that the generated means is meaningful and, where claimed, executable or admissible in the source practice. Then derive a new warranted consequence from its use. Selecting the next member of a fixed menu is a baseline, not the full target.

The construction rules for enlarging the repertoire remain explicit. A proposed intervention is distinguished from an implemented one; a syntactically valid rule from a sound one; a hypothetical law from acquired evidence. Establish a progress, capability, resource, or impossibility theorem for the expansion itself. This obligation links formation to the transformation of its own conditions without requiring every inquiry to run as one closed feedback algorithm.

\Ob{P17}{Make approximation, scale, and resources explanatory}
Extend a substantive exact formation or comparison theorem to an explicitly quantitative regime with an informative stability or approximation bound. Distinguish identifiable structure from stably recoverable structure, and separate statistical consistency from robustness of reconstruction. Explain how coarse graining, changed scale, idealization, or symmetry affects the intended answer and subsequent operation.

For supported compositions $W$, derive a law of the form
\[
 d_W\bigl(\Obs(WX),\Obs(W\widehat X)\bigr)
 \leq \omega_W(\varepsilon_1,\ldots,\varepsilon_m;\lambda),
\]
where the metrics, local errors, stability data $\lambda$, and domains are part of the theorem. The informative content must exceed a generic triangle inequality. Give a failure boundary for feedback, rare-event conditioning, long horizons, or scale loss when relevant, and a repaired positive regime rather than an unsupported uniform bound.

Distinguish costs of evidence, translation, proof or certificate verification, abstract construction, native implementation, and retaining a selected mechanism. Include correlation memory, clocks, shared randomness, communication, calibration, preprocessing, and oracle calls where used. A small carrier with an inaccessible decoder is not the same practical achievement as an executable adequate representation. Give exact or asymptotic resource laws and matching lower bounds where promised; do not require a polynomial-time solution to every realization problem.

\Ob{P18}{Explain formation in learning without making learning the subject}
Retain the response-defined learning-state target, and investigate a stronger formation result in which learning introduces or selects consequential representational organization. A sufficient state for a supplied learner is valuable but does not alone explain concept formation. Prove closure or controlled approximate closure, the relevant sufficiency or ambiguity property, and a consequence for future learning or use.

Derive a tractable structural presentation on a nontrivial family, rather than retaining all parameters, and establish a sharp boundary to stronger invariance or compression. The immediate cross-mechanism bridge for this local result follows the explicit relaxation in Appendix~\ref{app:ledger}; the programme's positive comparative obligation remains.

Compare against observable-algebra and nonlinear realization theory, Koopman-style methods, parametric learning, lenses, reverse differentiation, and actual representation-learning precedents. Distinguish coordinate transport from changing the optimizer, and behavior from training response. The elementary predictor $w=uv$ obeys $\dot w=-(u^2+v^2)\ell'(w)$ under its specified Euclidean gradient flow; equal present $w$ need not determine the next response. This remains a control, not a new formation theorem.

No neural benchmark is mandatory. If an empirical learning branch is selected, require prospective predictions fixed before the relevant future checkpoints or interventions are examined, held-out runs or interventions, and controls for architecture, initialization, objective, data order, and regularizer. A post hoc explanatory coordinate fitted to the outcome cannot satisfy a prospective discovery claim. A pure theorem can meet the relevant obligation without an empirical run, at its stated exact-input scope.

\section{The central integrated achievement}
\label{sec:integrated}

\Ob{P19}{Connect formation to warranted use and revision}
Produce an integrated positive result in which a constructed representation or theory supports a new warranted inference or action, survives the declared heterogeneous interactions, and admits a justified revision when its conditions change. The result must couple structural, deductive, epistemic, interventional, prescriptive, and computational content through actual dependencies. The components can be developed in different formalisms, and some links may be proved by separate theorems; their composition must itself be established.

This is a demanding positive target on a substantial class, not a compulsory pipeline through which every formation process must pass. Other results can explain transformations that lack recognition, deliberate selection, or a subsequent intervention.

\begin{target}[Formation to use to revision]
Choose a substantial class of independently specified, heterogeneous source presentations with partially overlapping meanings, explicit evidence or mathematical constraints, native operations, and a declared family of questions or tasks. Do not supply the sought shared carrier, language extension, joint realization, answer decoder, or compatibility certificate as the answer in disguise. Construct:
\begin{enumerate}
\item a consequential new structure, interface, or theoretical organization, with its dependence on the supplied material and its remaining alternatives;
\item an input-level adequacy theorem for an independently understood question family, including a warranted decoder or inference procedure using only available information;
\item a transport of the required warrants through actual contextual interactions, preserving the relevant interpretations, uncertainties, and operational preconditions;
\item a newly enabled inference, executable policy, or discriminating experiment that is nonvacuous and source-valid;
\item a typed comparison or revision after execution, new evidence, or language change, identifying conclusions retained, weakened, retracted, or newly justified;
\item an algorithm, constructive method, or exact characterization of the information and resources required for the claimed effective part, with a meaningful lower bound or obstruction;
\item a positive repair on a nontrivial class when one compatibility condition fails, retaining a useful conclusion instead of declaring permanent abstention.
\end{enumerate}
Prove the links between these outputs. A juxtaposition of separately correct algorithms or the identity representation followed by an always-abstaining policy does not satisfy the target.
\end{target}

\subsection{The quantifiers in a warranted action}

Let $\mathcal U(E)$ be a justified, nonempty family compatible with evidence $E$. When languages differ, the task comparison needed to interpret this family must first be supplied. Distinguish
\[
 \forall M\in\mathcal U(E)\ \exists a\in A(E)\;G(M,a)
 \qquad\text{from}\qquad
 \exists a\in A(E)\ \forall M\in\mathcal U(E)\;G(M,a).
\]
Here $A(E)$ is the independently specified set of source-executable actions available on the shared information. The first statement may use inaccessible knowledge of the actual model; it does not justify one robust shared action. Adaptive policies belong to a separately defined executable class $\Pi(E)$, may exploit future observations, and must be well-defined on the histories actually available. Their admissibility cannot secretly depend on knowing the actual model.

A probabilistic policy theorem might establish
\[
 \sup_{M\in\mathcal U(E)}
 \Pr^{\pi(E)}_M\{\text{declared task failure}\}\leq\varepsilon.
\]
Its probability semantics, horizon, conditioning, and uniformity must be specified. A separate confidence claim that the relevant truth belongs to $\mathcal U(E)$ with probability at least $1-\delta$ is a statement about the evidence procedure, not a probability distribution over worlds. A total failure bound such as $\delta+\varepsilon$ requires a proved link between these guarantees on the actual joint data-and-execution process. Adaptive selection, reuse of evidence, changed systems, and changed targets cannot be ignored. An empty uncertainty family indicates failure of assumptions or retained commitments, not a vacuous certificate of safety.

The deductive counterpart must similarly distinguish a proof in a representation from a proof that its interpreted conclusion applies to the source. The interventional counterpart must distinguish predicted effects from actually available controls. The computational counterpart must include the cost of acquiring and verifying the inputs to the policy. This is where the six dimensions meet mathematically.

Conditional guarantees require their conditioning information and domain. Laws do not determine unique conditionals on null histories; regular conditional kernels, when they exist, are generally determined only almost surely. A history with zero probability under one preparation may remain possible or become reachable after an intervention. Additional kernel, possibility, or access information must have its own warrant. Rare-event robustness requires appropriate positivity or regularity; the finite O7 repair below is not a theorem solving arbitrary measure-theoretic conditioning.

\subsection{Three connected flagship investigations}

\paragraph{Heterogeneous coordination.}
Construct a shared or interoperating account from agents with different observations, languages, dynamics, scales, and capabilities. Derive a consequential shared distinction and a valid collective inference or action without assuming one complete world model. Diagnose separately conflicting evidence, stale evidence, unmatched meanings, different viewpoints, model inadequacy, and real environmental change. Prove how a warranted revision propagates to dependent uses.

The route-passability control is exact: two marginal probabilities $\Pr(A)=\Pr(B)=1/2$ permit $\Pr(A\cup B)$ anywhere in $[1/2,1]$. A policy requiring a lower bound of $3/4$ is not robustly licensed by those marginals alone. The positive task is to determine which additional joint evidence, disclosed dependence assumption, or weaker action makes coordination warranted, and construct it on a nontrivial family. Data fusion alone is insufficient.

\paragraph{Scientific inquiry with changing theory.}
In a heterogeneous laboratory or scientific modelling network, an anomaly can require a new mechanism, measurement interpretation, or operation, not merely a parameter change. Connect P7 to actual experimental capability, uncertainty, deduction, and experiment choice. Derive a prospective consequence, identify the discriminator that could test it, and prove what survives when the revised theory is composed with other accounts. A laboratory twin is not successful merely because its simulation remains synchronized.

\paragraph{Mathematical formation with deductive warrant.}
Connect P8 to an account of concept introduction, proof transport, and revision under counterexamples. A common construction or new universal property must organize different originating theories and enable a substantive consequence. Compare this formation with the scientific case at the level of the principle that genuinely transfers. The comparison must preserve their different warrants; it may establish a directed bridge or a precise separation rather than one universal inference operation.

These are concrete programme-level ambitions, not an obligation to build robots, run a chemistry laboratory, or produce a general autonomous mathematician. Mathematical theorem families can establish them at a declared scope. The environmental or ecological case supplies an additional demanding contrast: spatial averaging that is adequate for flow can lose the local concentrations needed by an ecological interaction. Strategic or adversarial settings can test dependence on information selection and adaptive opponents where relevant, with objectives and authority to act kept distinct from predictive success. None of these diagnostics defines the subject or substitutes for its conceptual principles.

\section{The retained native construction commitments}
\label{sec:native}

The earlier native branch remains a demanding strand of the programme. It is not the default agenda for every formation result, and its technical detail is not a claim of foundational priority. It is also not erased by the broader formulation. The following obligations preserve its positive content; merely giving it a new interpretation does not complete it.

\Ob{P20}{Construct an unsupplied reusable interface}
From response patterns and a finite presentation of admissible contextual transformations, construct or obstruct a reusable factor presentation without taking the desired components, projections, or assembly as input. In the first finite test, exclude identity factorizations and one-part decompositions. Require at least two proper observable parts carrying distinctions used by actual tests or contexts, with an assembly recovering the claimed joint responses and operation action. Overlapping interfaces require a comparably substantive nontriviality criterion.

The output may be an identified family rather than a canonical decomposition. Positivity, correlated observations, predictive versus hidden state, and nonunique latent explanations must remain visible in the stochastic branch. A declared decomposition grammar is allowed; a supplied split and full joint witness is not an unsupplied-interface result. Extending finite verification to untested protocols requires a stated completeness argument.

An obstruction can settle a negative case and guide repair, but does not discharge this inherited positive construction commitment. A nontrivial positive family must remain or the target must be explicitly amended.

\Ob{P21}{Construct multistage native use and independent assembly}
For a nontrivial class of finite controlled systems, retain joint output and next-state kernels, preparations, partial local observations, specified retained models, and directed deterministic or stochastic comparisons. Require state reuse through at least two successive nonterminal transitions, genuine stochastic intermediate links, unknown joint choices, and a mechanism intervention or its precisely diagnosed failure. If feedback is claimed, allow a later action to depend on an earlier outcome.

For row-kernel conventions, write $M^i_{a,y}:S_i\to\RR_{\geq0}^{S_i}$ with $\sum_y M^i_{a,y}\one=\one$. A candidate joint implementation has carrier $V$, joint output and next-state kernels $X_{a,z}$, preparation $\pi_X$, and retained links $L_i$. For an observation map $q_i$, the required equation is
\[
 \left(\sum_{z:q_i(z)=y}X_{a,z}\right)L_i
    =L_i M^i_{a,y},\qquad \pi_XL_i=\pi_i,
\]
where $L_i:V\to\Delta S_i$ is the required stochastic link. Weighted sums replace deterministic output hiding for a stochastic observation channel. Corresponding equations must hold for the actual interventions. These statewise commitments are stronger than one initial trace law.

Derive an independent finite criterion that constructs joint native witnesses with no assumed global coupling or unknown global readout; retains the promised links or exactly identifies the weaker public behavior available; and compares at least two independently built legal assembly routes by actual native maps or a justified typed replacement. Rebracketing one supplied tensor is only a regression control. Prove a nontrivial native resource upper bound and a matching lower bound on a parametrized family, charging clocks, correlation carriers, and history. Explain how the finite observational certificate changes when a test, operation, or intervention is added or hidden.

The admissible class can be restricted by finite continuation bases, overlap structure, cone conditions, interface width, or stability. The restriction must not reduce every case to an already solved row polytope or known fixed-tensor gluing theorem. Reduce the inherited part to marginal extension, conditional products, supported coupling, semipullbacks, multicausal gluing, and bounded real-algebraic synthesis before isolating the new claim. A failed criterion requires a transparent counterexample and a positive repaired class; without the latter this positive milestone remains incomplete.

\Ob{P22}{Derive finite sufficiency with the promised native lifts}
Let $E\subseteq E^+$ be finite observation interfaces with restriction $q$, generating operations and exposed interventions, and native budget cap $B$. Let $H$ be an exact full response specification and $C$ the retained commitments. Write $\Real_b^+(H,C)$ for implementations of cost at most $b$ in the declared native class over $E^+$, and $\Real_b(qH,qC)$ for the restricted problem with correctly translated commitments $qC$. Derive finite input certificates, not assumed reverse functors, such that the required $E^+$ predictions are determined throughout generated protocols and, for every $b\leq B$,
\[
 \Real_b^+(H,C)\ne\varnothing
 \quad\Longleftrightarrow\quad
 \Real_b(qH,qC)\ne\varnothing.
\]
For the stated class of coarse implementations, construct same-budget native lifts and comparisons with the required composition law under admitted wiring and intervention operations. State whether the claim covers every implementation, specified endpoints, or only existence. The lifting conclusion is stronger than equality of feasible-budget profiles. Agreement only below $B$ does not determine minima above $B$.

Also construct a weaker useful certificate giving a directed comparison, controlled resource gap, or explicit missing-interface requirement when exactness fails. Derive the relevant full admissible profile information from available observational and operational inputs under explicit promises rather than supplying the sought native solution. Include a useful case with a predictive object larger than its hidden implementation and a case exposing failure of local adequacy after composition. Establish the corresponding algebraic relation-preservation instance without identifying its native category with the stochastic one.

The supplied-polytope condition $q(K)=K'$ with $\ker q\cap\operatorname{span}(K-K)=\{0\}$ remains an elementary control when full profile polytopes and joint correlation carriers are given. It does not solve the hard acquisition and native-lifting problem. Likewise a gap between generated and test-positive cones records freedom, and neither local cone equality nor marginal bases determines arbitrary joint correlations.

\Ob{P23}{Dovetail native results with cylinder normalization}
For every result claiming the supported native-to-cylinder comparison, construct an interpretation $J_{\mathfrak D}:\Real(\mathfrak D)\to\mathsf{CylReal}(\mathfrak D)$ of actual preparations, tests, native comparisons, and operations. State the information forgotten. Relate its native output to the established normalization span
\[
 X\xleftarrow{\mathcal M}G_X\xrightarrow{\mathcal E}R_{\mathfrak D}
\]
and a supported change to the established cross-doctrine comparison
\[
 F(R_{\mathfrak D})\xleftarrow{\mathcal M}G_F
                         \xrightarrow{\mathcal E}R_{\mathfrak D'}.
\]
The actual classes and variance are inherited from the source theorem. Prove the native/cylinder comparison by a typed map, two-cell, or span-between-apices, with units, associativity, and the existing compositor. For simultaneous signed changes use the surviving span construction where a direct map fails. Treat fixed-port wiring, changed wire types with reversed dual variance, and trace/feedback separately. Preservation of factorization classes does not by itself preserve the required mixed pullback or trace.

The new positive content must include reflection toward native semantics: a lifting criterion, resource cost, or exact obstruction when a formal normalized object lacks the promised native preimage. This obligation is activated by the supported comparison claim. It is not imposed on a pure logical formation theorem for which no cylinder interpretation is asserted. The clarification of this trigger is an explicit amendment to the earlier ambiguous universal imperative, not a waiver after making a native claim.

\section{The context and content bridge}
\label{sec:contextad}

Actual contextads and contentads are substantive mathematical candidates, not a synonym for a variable named context. Their grade objects, bodies, maps, substitution, transport, and coherence must be instantiated at the source problem's meanings. They are neither compulsory universal primitives nor decorative vocabulary.

\subsection{The achievements already available}

The relative-construction foundation uses a pure base $C$, category-valued modules, and the pseudocoend tensor $\odot_C$ in execution order. Its context and content modules admit a comma normal form retaining actual grade maps, noninvertible native cells, and changing-endpoint squares. A supplied coherent crossing gives mixed composition. This supplied-law route requires the crossing as additional input; it makes a weaker derivation claim than constructing the crossing from independently stated one-sided hypotheses. No general ordering of the two input regimes is asserted.

\paragraph{O1 independent interaction.}
O1 accepts a pullback- and composition-stable class $J$ of set maps containing identities and isomorphisms, a finitary cartesian polynomial monad preserving it, locally small coherent symmetric rig resources with finite augmented witness families, and an independently valid occurrence contentad. Its hypotheses include global opcartesian transport, genuine $J$-cartesian lifts, a natural polynomial occurrence map with its unit and accumulation laws, relative display pullbacks, preservation of cartesian arrows, and invertibility of the specified canonical Beck--Chevalley mates. It constructs the crossing, four invertible modifications, and full coherence on the original category-valued modules. It does not merely assume a distributive law. Interface sets may be arbitrary, while witness fibres and polynomial positions are finite.

Its necessity result is confined to the marked exhaustive occurrence class and realizable projection maps $D$. An independent richness argument $D=J$ gives necessity and sufficiency only within that same marked class, not a classification of unrestricted interaction laws. Strong changes of independent inputs already have their own coherent comparison theorem. Neither a smooth parameter manifold nor arbitrary infinitary content is covered merely by being called a parameter. These are achieved scoped dependencies, not claims newly proved here \cite{O1,Relative}.

\paragraph{O5 contextual replacement.}
Finite primitive state-and-cost couplings, a full-state relation, and a compatible initial coupling yield replacement for every admitted finite typed client, including adaptive branches and repeated method use. The proof constructs paired-history native spans. The client class is not a finite suite: finite terms can have arbitrarily large size within the admitted finite ranked regime. Observations retain the joint law of cost and public causal tree, including ordered leaf positions. Every fixed client primitive must preserve the relation; identical names are not enough. This is sufficient contextual soundness, not a complete characterization of contextual refinement. Rational local coupling feasibility is a finite linear problem, whereas arbitrary real coefficients do not automatically give a decision procedure. General concurrency, unrestricted recursion, and physical recovery are not consequences. Recovery requires the stated capabilities and residual class \cite{O5}.

\paragraph{O7 law and possibility.}
O7 repairs the loss of zero-probability possibilities by retaining both a law and a possibility set. Writing $\mathcal P_f(X)$ for finite-support probability laws and $\mathcal P_{\mathrm{fin}}(X)$ for finite subsets, its construction uses
\[
 \begin{gathered}
 \mathsf{M}X=\{(\mu,S)\in\mathcal P_f(X)\times\mathcal P_{\mathrm{fin}}(X):
                         \supp\mu\subseteq S\},\\
 \mathsf{K}X=\mathsf{M}(\mathsf{C}\times\mathsf{T}X),
 \end{gathered}
\]
with the specified costs and typed occurrence-tree outcomes. Equality captures the joint law and all universal predicates on retained possible outcomes, not raw witness multiplicities or all parallel cells. A finite law coupling plus candidate possibility coverage is equivalent to a finite original-O1 paired-history span with surjective candidate-history projection, cost-exact reference leg, cost-lowering candidate leg, and the specified public-body/private-state compatibility. Coverage is not positive support: a zero-weight candidate possibility still needs an appropriate reference possibility.

The presence-enriched repair is itself an independent O1 input. All raw objects lift as all-present objects, but exactly surjective old native maps lift between those all-present objects; the repair does not preserve every old cell. The reconstruction is a fixed-endpoint reflection with restricted cost-exact native maps and generally noninvertible lax comparison from reconstructed composites to reconstruction of a composite. It is not a global equivalence, and its ambient representatives need not be definable by the source primitives or physically implementable. Even an exact probability-law oracle does not recover independent zero-weight possibility data. This is a limitation of the chosen semantics, not a failure of O1's interaction theorem \cite{O7}.

\Ob{P24}{Construct the missing phenomenological bridge}
Using these results at full strength where applicable, establish a new positive formation, acquisition, language-change, or native-lifting theorem connecting independently warranted source presentations to usable operational conclusions. The preferred difficult interfaces are: generating adequate shared interfaces and their certificates rather than assuming them; acquiring the probability and possibility information actually required by replacement; constructing a revised interface after a new operation or interpretation; and characterizing source-definable lifts of an adequate ambient semantics.

A selected theorem must construct its observational or answer-level operational structure and prove closure before writing a pseudopromonad for it. Specify category-valued modules $H^{\mathrm n}:C^{\mathrm{op}}\times C\to\Cat$ and $H^{\mathrm o}:\bar C^{\mathrm{op}}\times\bar C\to\Cat$, with justified size or nested-universe assumptions for the pseudocoends. Their realizations have designated pure pseudofunctors $J^{\mathrm n}:C\to B_{H^{\mathrm n}}$ and $J^{\mathrm o}:\bar C\to B_{H^{\mathrm o}}$. If a strong comparison is claimed, provide a functor $F$ and pseudofunctor $q$ with
\[
 \begin{gathered}
 F:C\longrightarrow\bar C,\qquad
 q:B_{H^{\mathrm n}}\longrightarrow B_{H^{\mathrm o}},\\
 q(a)=F(a),\qquad qJ^{\mathrm n}\xRightarrow{\sim}J^{\mathrm o}F,
 \end{gathered}
\]
invertible unit and composition comparisons and the pure-compatible invertible icon displayed above. A lax or oplax result must instead display its actual compositor direction and prove only the conclusions it supports. The native module has unit $I_C\to H^{\mathrm n}$ and multiplication $H^{\mathrm n}\odot_C H^{\mathrm n}\to H^{\mathrm n}$, where $I_C(a,b)=\operatorname{disc}C(a,b)$; the observational module has its separate unit $I_{\bar C}\to H^{\mathrm o}$ and multiplication over $\bar C$. Both require their own coherence. A computation in a hom category is not automatically a whole native realization.

O7's fixed-endpoint cost-exact reflection does not by itself provide this comparison on the full changing-endpoint, cost-lowering native calculus. The needed target comparison categories or appropriate restrictions must be constructed separately.

Distinguish three further achievements: an ambient native witness; a witness definable in the permitted source grammar with its resources and ownership; and an executable physical implementation. A forward operational comparison supplies none of their inverses. Use witness categories or typed spans when coherent selection fails. When the cylinder claim is made, connect the model-level lift to P23 rather than substituting an operational notation for it.

Any claimed finite or effective criterion must yield more than the elementary statement that target probes pull back to admitted source probes. Identify the residual after O1's derived crossing, O5's client induction, O7's repair, existing language-change results, and native/cylinder normalization have been used. Generic finite grammar enumeration is not automatically that residual. A valid result must construct an unsupplied adequate object, derive a new acquisition or native-resource law, or establish a new interaction between formation and contextual use, with a sharp boundary and positive surviving class.

The route remains a candidate. If it cannot faithfully express the selected question, use another formalism and state the non-applicability precisely. A comparison study of its strengths and limits is owed; forcing all deductive, empirical, or reflective formation into it is not.

\section{Generality and the programme level result}
\label{sec:success}

\Ob{P25}{Prove common mathematics across genuinely different settings}
The programme must obtain common laws whose actual hypotheses and conclusions do work across different evidence, operation, and answer semantics. Retain the inherited requirement to instantiate the same substantive theorem in algebraic and stochastic/positive settings, and work through a finite-dimensional linear setting or state precisely which additional analytic assumptions it needs. Static and sequential cases must be represented in the advertised common account. These named tests are a retained obligation, not the outer boundary of the subject.

Extend the comparison to mathematical and scientific theory formation with their different warrants. The common content may be restricted, directed, or mediated by an interoperating family of doctrines. It must produce a consequence in the receiving setting, not simply provide vocabulary. Failure of a proposed transport must identify the assumption or preserved question responsible. Prove any claim of unavoidable heterogeneity over a specified class of reductions, rather than infer it from several unsuccessful encodings.

Unfamiliar cases should enter while they can still change a principle. A theorem need not cover every domain or earn immediate cross-domain relevance to be mathematically successful; the programme nevertheless cannot declare generality while all substantive results remain variants of one finite architecture.

\Ob{P26}{Deliver an explanatory synthesis with new consequences}
Produce a coherent mathematical account that answers what has become understood about formation, warrant, heterogeneous use, and revision. Identify a question left unresolved by the separate results and settle it by combining them. The synthesis must yield a construction, theorem, operational procedure, discriminating consequence, or exact classification that is not supplied by merely placing the components side by side.

Its central claims must be matched to actual results: what accounts can be formed; why their contents are warranted; how they acquire additional expressive or operational power; when they compose; what forces revision; and why their stability is more than persistence of an unchallenged convention. A synthesis of known mathematics may contribute substantially, but the independent promises of new positive results in P5 and P6 remain. Neither theorem novelty nor the number of manuscripts measures the achievement on its own.

\subsection{What counts as the ambitious outcome}

The programme earns its central claim through a connected body of results, not necessarily one universal theorem. That body must establish:
\begin{enumerate}
\item a discriminating account of the subject with mathematically consequential boundaries;
\item positive construction and positive comparison beyond the established baselines;
\item scientific and mathematical formation results that cannot be substituted by fixed-language state estimation;
\item a warranted integration of unlike explanations, including a nontrivial source-valid use and a justified response to change;
\item an account of immanent transformation, reflective understanding where claimed, and a proved form of stabilization or persistence;
\item genuine cross-setting consequences, informative quantitative or effective fragments, and exact limits.
\end{enumerate}
The inherited positive native commitments remain additional named obligations of the proposal. Completing another strand cannot mark them fulfilled. A full-contract completion claim requires their completion or an explicit accepted amendment recording their replacement or retirement. The relative-computation route itself remains conditional, with the separate achieved programme untouched.

There is no fixed ceiling on new obligations required by the mathematics. A discovered dependence can enlarge the programme; a refuted organizing hypothesis can redirect it. Neither change authorizes retroactive alteration of a completed local theorem's acceptance conditions. Ambition is enlarged by consequential questions, not by forcing every project to carry every obligation.

\section{How the work is selected and judged}
\label{sec:practice}

This contract sets intellectual commitments, not a rigid sequence of gates. Construction, comparison, formation, and integration should constrain one another from the start. Evidence and operational access shape an empirical problem before its exact model is frozen. A mathematical formation problem begins with its own judgments, consequences, and interpretations rather than an artificial empirical proxy. Heterogeneous tests can challenge a candidate principle before it is polished into an all-purpose formalism.

Before a substantial investigation, state its question, input class, answer requirements, independent hypotheses, promised positive or negative conclusion, strongest relevant predecessor, and decisive failure condition. Give the actual local theorem or constructive objective rather than treating this programme contract as if it had already selected every theorem's assumptions. Choose the investigation for the consequential uncertainty it can resolve about the programme, not because it is the nearest extension of the previous calculation.

At a major change of direction, make one substantive judgment: what is now understood that was not understood before, which competing accounts remain, and what next question would materially change the answer? This is an intellectual selection rule, not a demand for a new audit or report after every lemma. A valid specialized theorem can remain useful supporting mathematics while ceasing to determine the agenda.

\paragraph{Four judgments are separate.}
Mathematical correctness asks whether the stated result follows. Explanatory contribution asks what it clarifies. Programme significance asks which central question it advances. Novelty asks what is new relative to actual predecessors. None entails the others. A synthesis can be explanatory without priority; a new difficult theorem can be peripheral; an obstruction can redirect the subject without fulfilling a positive construction promise.

\paragraph{Proof standard.}
For each substantive claim give a complete statement and proof with typed maps and diagrams, all size, positivity, regularity, reachability, interpretation, access, and resource assumptions, and explicit supplied/derived boundaries. Keep native object, operation, comparison, and cost meanings visible. Use exact examples and reproducible finite checks where they test the claimed formulas, including empty fibres, zero masses, unused states, undefined operations, and infinite-cost cases. Require independent mathematical review and separate source/provenance review. Computation and review are evidence at their actual scope, not proof-assistant verification or historical novelty clearance.

\paragraph{Computational standard.}
Separate effective enumeration, decidability, witness verification, polynomial or other complexity, approximation guarantees, experimental feasibility, and heuristics. A stationary point, small residual, global zero, and certified positive global minimum are different results. State bit complexity for exact algorithmic claims and report numerical failures honestly. For extended nonnegative real costs under a lawful witness map $t$, with $\inf\varnothing=+\infty$, the familiar inherited cost
\[
 (t_!c)(z')=\inf\{c(z):t(z)=z'\}
\]
has an associative elimination law; that does not make it equal to the independently defined cost of a target native problem. Zero sets, residual norms, gradients, and optimizer trajectories must each have their own transport claim. No surrogate may replace an inconvenient promised objective without declaring a change of problem.

\paragraph{Stopping and repair.}
Stop asserting a local theorem when a verified counterexample defeats it, the whole claim reduces to prior work, or all surviving premises presuppose the sought result. Numerical failure is not an impossibility proof. An unavailable experiment blocks the empirical conclusion; it does not license an oracle. Record proved as stated, open with the exact gap, or refuted as stated. A repair states its changed semantics and proves a useful positive class. A restricted theorem retains its scope and credit. Reprioritization, replacement, and retirement are distinct from completion.

\paragraph{Deliverables.}
The outputs are a body of proved formation and comparison results, the integrated result of P19, source-faithful constructions and counterexamples, exact or reproducible computational witnesses where claimed, and the explanatory synthesis of P26. Their packaging may be manuscripts, formal developments, implementations, or worked mathematical investigations as appropriate. There is no manuscript quota, compulsory software stack, publication promise, implementation mandate, or execution schedule in this proposal.

\section{Decisions left for adoption}
\label{sec:adoption}

This proposal makes substantive choices rather than merely asking what a contract could contain: formation of both mathematical and scientific theories is central; heterogeneity is a candidate constitutive feature; recognition is not required for every transformation; and integration is an independent positive achievement. The following questions would most usefully sharpen adoption, without blocking the present proposal.

\begin{enumerate}
\item Does the stated intellectual prize capture the desired centre, especially formation of questions, languages, and warrants rather than only reconstruction and reuse?
\item Should the first bounded investigation challenge the subject's formulation, construct a theory-formation principle, or attempt the integrated formation-to-use theorem? No answer is presumed by the order of the sections.
\item Are the inherited native flagship and named comparative trio to remain explicit long-range commitments as proposed, or should a particular target be deliberately replaced? They have not been retired here.
\item Is the proposed relaxation of immediate cross-mechanism comparison for every local generative theorem acceptable, with the programme's positive bridge and integration burden retained?
\item Is a well-characterized interoperating family, with proved limits to stronger unification, an adequate architectural outcome when all its positive formation and use commitments are met? This proposal says yes, while leaving unity open to the mathematics.
\end{enumerate}

Acceptance of the document would adopt a research contract, not prove its hypotheses. The present draft does not itself replace an accepted agreement, reopen the separate relative-computation contract, or begin an unrestricted research schedule. It is intentionally uncompiled \LaTeX{} source; only source-level checks and review are claimed.

\appendix
\section{Carry forward and deliberate change ledger}
\label{app:ledger}

The reference point is the actual amended \emph{A research contract for categorical phenomenology}, dated 3 November 2026, not a caricature of it as a finite-model programme \cite{Earlier}. It already protected broad generation, comparison, alternative formalisms, and conditional native branches. The reassessment diagnosed an asymmetry of operative detail and research selection, not wholesale absence of these principles \cite{Reassessment}. The dispositions below apply only if this proposal is adopted.

\begin{description}[style=nextline,leftmargin=0pt]
\item[Subject \quad broadened]
The observation-to-structure orientation is retained as a substantial strand, not the definition of the whole subject. P1--P4 add discriminating subject formulation, question formation, typed warrant, and obtainable access. The new scope includes judgments, mathematical constraints, theory languages, and immanent transformation without pretending they all share one observational primitive.

\item[Paired positive commitments \quad carried and enlarged]
The independent positive generative and comparative targets survive in P5 and P6. New positive formation and integration commitments are added in P7--P19 and P26. A negative classification or more downstream papers cannot compensate for an unfulfilled positive target.

\item[The six generative nonvacuity requirements \quad carried]
Specified output/equivalence/use; visible supplied structure; independent input criterion; positive nontrivial family; counterfactual dependence and relevant invariance; warrant, ambiguity, and predecessor comparison remain binding under P5. Abstract noncircular assumptions are permitted; effectiveness is an additional proved claim, and the finite native branches retain their finite-criterion requirement.

\item[Gate A \quad rigor retained and selection broadened]
The earlier requirement was to state the actual theorem, semantics, input class, positive conclusion, independent assumptions, and a residual beyond the strongest predecessor, and to compare organizing languages. Those requirements remain for a bounded theorem project. The programme contract itself is not falsely called the already selected first theorem. The original initial comparison with incidence/FCA, predictive reconstruction, CALF, realization, Stacey, and cylinder work remains required for the inherited comparative branch. A different foundational or formation project uses the strongest sources for its actual claim; this is an explicit change of the universal opening source roster, not permission to evade relevant predecessors.

\item[Gate B \quad one explicit relaxation]
The original generative target required a substantive relation to another reconstruction mechanism inside each accepted generative result. This proposal relocates that local requirement to the cumulative positive comparison and integration obligations P6, P19, P25, and P26. A local formation theorem can succeed on its own stated terms. It does not thereby pass the old unchanged gate, and the missing programme bridge remains open. The positive theorem, nontrivial-family, boundary, and six nonvacuity floors are not relaxed.

\item[Gate C and the named trio \quad carried and made nonexclusive]
Incidence/closure, temporal/predictive, and native-factorization mechanisms must still receive the promised positive comparison at their correct types, beyond generated-image and cylinder baselines. P6 retains this open target. Additional contrasts can now be equally central; they do not silently complete or replace the native-factorization portion. Directed and partial comparisons remain legitimate, with exact loss and impossibility scope.

\item[Gate D and the native flagship \quad carried with explicit triggers]
P20--P22 preserve the unsupplied-interface, independent multistage assembly, stochastic links, sequential reuse, mechanism intervention, same-budget lifting, and matching resource-bound demands. Other formation results need not complete them before receiving local credit. This was already permitted in the old text; making the trigger visible does not weaken their conclusions. Retirement or reduction would require a separate explicit amendment.

\item[Cylinder and relative-computation bridge \quad trigger clarified]
P23 requires the actual native/model-level normalization comparison whenever that supported route is claimed. It explicitly narrows the old ambiguous phrase ``The new theory must'' so that an unrelated logical theorem does not acquire an unclaimed cylinder burden. Once a comparison is claimed, its variance, compositors, span-between-apices, joint wiring, native reflection, and trace restrictions remain intact. P24 retains the conditional relative-computation route and strengthens its new-result ambition; O1/O5/O7 remain achieved dependencies, not renamed obligations.

\item[Gate E \quad carried and moved into problem design]
Exact laws, exact queries, samples, destructive sequential access, and hypothetical extensions remain separate. An empirical claim must establish executable access, justified uncertainty or convergence, and resource costs; exact mathematics does not inherit a compulsory experiment. Access now helps shape the question before a preferred model is fixed. Deductive theory formation has its own warrant rather than an empirical gate by analogy.

\item[Gate F and generality \quad carried and assessed throughout]
P25 retains the same-substantive-law requirement across algebraic, stochastic/positive, and appropriately qualified finite-dimensional linear cases, including static/sequential contrasts. P26 adds a consequential synthesis at programme level. Major research choices use explanatory significance throughout, rather than reserving it for the final gate. Coherence among papers is not a substitute for the promised common theorem or translated proof.

\item[Learning \quad retained as a candidate and strengthened]
P18 retains the response-defined sufficient-state target and its scalar/gauge, architecture, access, source, and prospective empirical controls. It adds a stronger formation ambition without claiming the old target already achieved it. A neural experiment remains optional, and the programme is not defined by a supplied optimizer.

\item[Resources and proof discipline \quad carried]
The distinct existence, endpoint-comparison, universal-lift, budget, uncertainty, and native-resource quantifiers remain. Exact verification, full proofs, typed comparisons, meaningful edge cases, mathematical review, source review, and an honest result-status ledger remain required. Inherited infimization, global optimization, and gradient assembly are not conflated. No novelty clearance follows from review alone.

\item[Stress tests \quad content carried, completeness role revised]
The 21 mathematical controls in Appendix~\ref{app:controls} remain mandatory when their semantics are claimed. The old requirement that whole-programme generality eventually cover the full suite is retained for claims spanning those semantics, but is not a completeness specification for the broadened subject: new formulations can require different tests. Passing the suite does not discharge P1 or establish universal scope. Replacing a redundant example must preserve the failure mode it diagnoses.

\item[Research order \quad revised]
The old diagram placed native, empirical, and cross-instance tests after initial comparative/generative results. This proposal replaces that order with interacting investigations and early challenge by relevant access and heterogeneity. There is no gate for every lemma and no compulsory path through one ambient technology. Local achievement remains judged against its actual stated conditions.
\end{description}

\section{Inherited controls and added discriminators}
\label{app:controls}

For each substantial investigation, record which controls below apply to its claims and which lie outside its scope; a reason can cover a whole inapplicable group and need not be repeated for every lemma. A negative answer can be correct. It cannot be concealed by silently changing the native class, interpretation, or requested use. These are diagnostic mathematics, not a checklist that defines the subject.

\begin{enumerate}
\item \textbf{Hidden output in the input.} Change the doctrine or grammar and track which generated structure changes. A fixed pairing can support meet completion, linear span, or native factorization with different answers.
\item \textbf{Evidence versus access versus law.} Separate more records, newly available tests, and a hypothetical larger law. Confidence families need not be nested unless their construction proves this.
\item \textbf{Parameters versus observational bridge.} Check reparameterization direction, operation descent, native/cylinder typing, finite witness hypotheses, cell retention, and the direction of lax reconstruction where invoked.
\item \textbf{Present equality versus learning response.} Retain the $w=uv$ counterexample, including stationary-loss and restricted-probe limits. A sufficient learning state must close under its declared protocols.
\item \textbf{Coordinates versus a new system.} Transport dynamics and readout together. Imposing a fresh Euclidean optimizer after a nonlinear coordinate change is not a gauge transformation.
\item \textbf{Preinstalled decomposition.} Architecture, regularizer, initialization, and vocabulary can determine the supposed discovery. Separate identification, selection, and fitting.
\item \textbf{Non-equivalence versus non-comparison.} Failure of an inverse, adjoint, or proposed functor does not exclude directed maps, spans, or another interpretation.
\item \textbf{Observations versus operations.} Separating response profiles do not supply a multiplication absent a composition doctrine or composite observations.
\item \textbf{Noncanonical emergence.} Compatible latent factors can remain inequivalent. A prior or compression score selects without necessarily identifying.
\item \textbf{Row separation versus completion.} Rows $10$ and $01$ remain separated after keeping the first coordinate, but their meet/top completion loses distinctions. Preserve the generated relations claimed, not just original rows.
\item \textbf{Abstract hull versus native geometry.} Affinely isomorphic response squares can have different native minimum costs in different embeddings. Retain the admissible profile class.
\item \textbf{Predictive versus hidden size.} Finite hidden models can have infinitely many predictive states. Hiding an output may require refiltering rather than a map on old states.
\item \textbf{Local versus joint evidence.} Rank-one local behaviors can have rank-two joint behavior under correlated preparations. Marginal reconstruction does not create joint information.
\item \textbf{Chosen mechanism versus refinement.} Equal coarse traces for early/late randomness need not preserve a specified statewise anchor after a signal is exposed. More states do not universally repair the conflict.
\item \textbf{Support versus quantitative feasibility.} Equal full-support patterns can have different stochastic feasibility and minimum sizes. O7's extra possibility information does not erase quantitative distinctions.
\item \textbf{One-step feasibility versus independent assembly.} Independently selected intermediate completions may lack a simultaneous cover. Rebracketing a supplied tensor does not decide this problem.
\item \textbf{Local versus joint memory.} One-state local factors can require unbounded joint correlation memory. Count shared randomness and phase clocks.
\item \textbf{Optimum value versus coherent choice.} Symmetry can forbid an equivariant minimum-state selector. A family or resource correspondence may be the correct output.
\item \textbf{Feasibility versus optimizer transport.} Test the actual quartic's zero sets, residuals, gradients, and values where that solver is claimed. A convenient surrogate is not its proof.
\item \textbf{Finite information versus future behavior.} Tested protocols can leave an untested continuation or intervention unresolved. State the class promise and available discriminator.
\item \textbf{Feedback versus local error.} Give stability assumptions, a meaningful bound, or a rigorous failure of uniform control.
\end{enumerate}

The broadened subject adds further decisive tests: multivalued adequacy without a coarsest quotient; a valid uncertainty region rather than a forced model choice; an expressive change versus a definitional or coordinate change; transformation without recognition and recognition without transformation; individually justified claims with an unjustified composite; order-dependent revision; useful coordination without a single global realization; and a mathematical formation problem whose warrant cannot be reduced to sampling. These tests may be replaced or sharpened when the theory explains why. Their relevance cannot be evaded by defining the problematic question out of scope while continuing to advertise its answer.

\section{The baseline and the unresolved advance}
\label{app:baseline}

The earlier contract's generated-image, predictive, positive-realization, cylinder, finite-native, and relative-computation results remain dependencies at their actual hypotheses and review status. This proposal has not independently reproved all of them. It uses the source-level reassessment and the inspected O1/O5/O7 statements to avoid treating achieved work as future novelty.

Later work has already supplied a meaningful operational bridge among residual automata, formal concepts, and topology with exact partial-action domains; distinct extension criteria for FCA, future behavior, and Dowker structure; structure-targeted and sequential evidence results; a conditional positive source-law identification with a further predicted moment consequence; and evidence-guided language revision beyond fixed linear capacity. These are real achievements. Their assumptions about candidate families, admitted access, exact laws, proposal grammars, and interpretation remain visible. A predicted held-out consequence is not newly collected empirical confirmation.

The remaining programme-level advance is not simply another instance of any one of these constructions. It is a coherent account of formation and warrant across changing languages and heterogeneous practices, together with new positive consequences and the source-level use/revision theorem. The named native-factorization comparison, unsupplied-interface multistage target, promised data-derived native sufficiency, and new relative-computation bridge are not declared completed by the preceding notes. Any later completion claim must match an actual theorem to each specific clause rather than inherit an indefinitely stale list of gaps or presume success from a related title.

\section{Precedents that set the level of the ambition}
\label{app:precedents}

The sources below constrain claims of novelty and indicate genuine mathematical resources. Their conjunction is not evidence that the proposed integration is absent from the literature. Targeted primary-source checks support the stated controls; they do not amount to independently reproving every source or reproducing its experiments. Before asserting a new theorem, narrow its exact hypotheses and conclusion and compare them with the strongest relevant results and successors.

\paragraph{Construction and comparative mathematics.}
Generated/observable factorization, CALF, classical realization, and Garner's cylinder-factorization and Isbell work already supply substantial construction theorems \cite{CALF,Garner}. Stacey's comparisons depend on particular test, forcing, and adequacy hypotheses; inclusions, adjunctions, and equivalences must not be inferred interchangeably \cite{Stacey}. Abstract interpretation and generalized strong preservation already connect observation languages, abstract domains, and complete-domain construction \cite{StrongPreservation}. A new object generated by a classical construction can have explanatory value, but its universal property is not new merely because it appears in phenomenology.

\paragraph{Question-relative information.}
For finite experiments over a common parameter set, dominance for every finite decision problem is equivalent to parameter-independent stochastic postprocessing. The finite Blackwell theorem was checked through the explicit formulation in Buscemi, not by claiming to have independently inspected Blackwell's original proof \cite{Blackwell,Buscemi}. Restricting the tasks gives a weaker comparison and does not entail full garbling or recoverability. Categorical statistical sufficiency has additional hypotheses, including strict positivity in the specific recovery characterization used as a control; a supplied statistical model is not a finite-sample identification result \cite{Fritz}. These are strong precedents for P2--P4, not evidence that all heterogeneous questions share one information order.

\paragraph{Heterogeneous logic and proof.}
Institutions and the heterogeneous-theory framework already organize signature change, satisfaction, theories, interpretations, conservativity, and heterogeneous proving. Mossakowski's construction retains component logics and relates them rather than demanding one combined logic \cite{Hets}. Model-theoretic, consequence, and proof-theoretic conservativity differ. A new graph of doctrines or satisfaction-preserving translations is therefore infrastructure, not the intended advance. The residual must concern formation, warrant across interpretations, operational use, or another consequence not already obtained there. The method map explains that its apparent earlier aversion to institutions was accidental; the preceding contract already allowed them. No such restriction is imposed here.

\paragraph{Inference, learning, and control already interact.}
Parametric optics and categorical cybernetics already relate learners, controllers, and games \cite{Cybernetics}; compositional gradient learning already combines parametrization, lenses, and reverse differentiation \cite{Gradient}. Exact Bayesian lens composition yields inversion of the composite up to the appropriate almost-equality and support qualifications, and approximate-inference doctrines connect specified statistical models to hierarchical dynamics \cite{SmitheI,SmitheII}. None of these statements by itself chooses a scientifically adequate language, validates its objective, or guarantees global convergence. Conversely, it would be inaccurate to present the predecessors as merely a collection of disconnected models.

\paragraph{Composition and its losses.}
Open Petri-net operational semantics preserves composition in its richer setting, while the reachability summary is only lax: repeated boundary crossing can create behavior missing from composed endpoint summaries \cite{Petri}. Causal abstraction already has explicit interventional error composition; for the finite setting and chosen square-root Jensen--Shannon discrepancy of Rischel and Weichwald, the abstraction error is subadditive under composition \cite{CausalError}. Neither result is a generic theorem for arbitrary summaries or metrics. Sheaf-theoretic contextuality provides exact local-to-global obstruction theorems with specified measurement identities, covers, and outcome laws, not a universal verdict on heterogeneous model collections \cite{Contextuality}. These are positive standards for P13--P17, not new obstructions to claim again.

\paragraph{Theory formation is existing territory.}
HR already introduced definitions through production rules, formed conjectures, used theorem provers and countermodel finders, and fed the results back into theory formation \cite{HR}. DreamCoder already learns reusable library abstractions and search structure within a supplied metalanguage \cite{DreamCoder}. Definitional or computational gains can be substantial without expanding denotational expressivity. Recent scientific-discovery preprints also study verified schema transitions with categorical transport and model-class revision through coupled edits and experiments \cite{ScientificSchema,Outgrow}. Their stated grammars, validation gates, predictive assumptions, and bounded proposal spaces remain supplied; their existence makes ``language change'' or ``categorical revision'' alone an inadequate novelty claim. These recent sources were inspected at their stated results, not independently reproduced.

The programme should use these precedents to raise its level: establish a consequential formation or integration theorem that they do not already provide under the same hypotheses, or give a rigorous synthesis whose new explanatory consequence is explicit. It should not search for novelty by renaming known mechanisms, nor discard a powerful predecessor merely because it is not written in the preferred vocabulary.

\section{Sources and status of this proposal}
\label{app:sources}

The user's complete conceptual formulation, the original method map, the actual preceding amended contract, and the reassessment supply the proposal's orientation and inherited obligations \cite{Formulation,Map,Earlier,Reassessment}. The contextual-realization problem note supplies a developed strand and earlier theorem-shaped tests, not a compulsory universal foundation \cite{Problem}. The relative-computation sources are used at the exact O1/O5/O7 scope recorded above. Later programme achievements are credited through the reassessment and their identified source notes; their completion and priority status are not inferred from filenames or retrieval dates.

This document states a research programme and its proposed acceptance obligations. It proves no new programme theorem and certifies no empirical finding. Static source checks cover balanced environments, cross-references, and bibliography keys. The source remains intentionally uncompiled, so no PDF layout or compilation is claimed. Adoption remains a separate decision.

\begin{thebibliography}{99}

\bibitem{Formulation}
\emph{Categorical Phenomenology: Subject, Mathematical Formulation, and Research Objectives.}
User-supplied conceptual document, including the four motivating problems. Complete formulation consulted; its interfaces and subject boundary are explicitly provisional.

\bibitem{Map}
\emph{Methods for generative categorical phenomenology.}
Original programme method map, with forty-two thematic sections and candidate entry points. Its literature leads and dependency sketch are non-exhaustive and provisional.

\bibitem{Problem}
\emph{Categorical Phenomenology: A Problem Note on Contextual Realization.}
Original programme research note. Observational realization, contextual saturation, decomposition, intervention, typing, stability, and transfer are a developed candidate strand.

\bibitem{Earlier}
\emph{A research contract for categorical phenomenology.}
Revised discussion draft with parametric amendment, dated 3 November 2026. The present proposal compares its actual targets, six gates, twenty-one controls, and source qualifications; its historical adoption status is not inferred from its own draft notice.

\bibitem{Reassessment}
\emph{Phenomenology contract reevaluation.}
Programme assessment, 4 November 2026. Includes the source-level treatment of the common operational construction, generated characterization, evidence and access, two-stage source-law identification, and characterization-language revision.

\bibitem{Relative}
\emph{One relative construction for context content and effects.}
Programme source, 1 November 2026. Relative pseudopromonads, category-valued native cells, changing-endpoint squares, comma normalization, and the supplied-law factory.

\bibitem{O1}
\emph{General Independent Interaction Criterion: Augmented rig contexts and relative occurrence displayed content.}
Programme proof dossier, 1 November 2026. Independent relative interaction, forced content exactness in the marked class, and change-of-input comparisons. Mathematical proof, not claimed proof-assistant certification.

\bibitem{O5}
\emph{Contextual Refinement and Controlled Recovery: Finite stateful replacement through retained probabilistic cells.}
Programme proof dossier, 1 November 2026. Local coupling criteria, open contextual soundness, native paired histories, and capability-bounded recovery.

\bibitem{O7}
\emph{Law and Possibility Repair for Dependent Protocols: A compositional finite repair of the all path observation gap.}
Programme proof dossier, 1 November 2026. Law-and-possibility semantics, covered native spans, presence enrichment, and local lax reflection.

\bibitem{CALF}
G. van Heerdt, M. Sammartino, and A. Silva.
\emph{CALF: Categorical Automata Learning Framework.}
CSL 2017.
\url{https://arxiv.org/abs/1704.05676}.

\bibitem{Garner}
R. Garner. \emph{The Isbell monad.}
Advances in Mathematics 274 (2015), 516--537.
\url{https://arxiv.org/abs/1410.7108}.

\bibitem{Stacey}
A. Stacey. \emph{Comparative Smootheology.}
\url{https://arxiv.org/abs/0802.2225}.

\bibitem{StrongPreservation}
F. Ranzato and F. Tapparo.
\emph{Generalized Strong Preservation by Abstract Interpretation.}
Journal of Logic and Computation, 2007.
\url{https://arxiv.org/abs/cs/0401016}.

\bibitem{Blackwell}
D. Blackwell. \emph{Equivalent Comparisons of Experiments.}
Annals of Mathematical Statistics 24(2) (1953), 265--272.
\url{https://doi.org/10.1214/aoms/1177729032}.

\bibitem{Buscemi}
F. Buscemi. \emph{Comparison of quantum statistical models: equivalent conditions for sufficiency.}
Communications in Mathematical Physics 310 (2012), 625--647. The finite classical control used here is the explicitly stated Blackwell theorem in Section 2.
\url{https://arxiv.org/abs/1004.3794}.

\bibitem{Fritz}
T. Fritz. \emph{A synthetic approach to Markov kernels, conditional independence and theorems on sufficient statistics.}
Advances in Mathematics 370 (2020), 107239.
\url{https://arxiv.org/abs/1908.07021}.

\bibitem{Hets}
T. Mossakowski. \emph{Heterogeneous Theories and the Heterogeneous Tool Set.}
Dagstuhl Seminar Proceedings 04391, 2005.
\url{https://doi.org/10.4230/DagSemProc.04391.7}.

\bibitem{Cybernetics}
M. Capucci, B. Gavranovi\'{c}, J. Hedges, and E. F. Rischel.
\emph{Towards Foundations of Categorical Cybernetics.}
EPTCS 372 (2022), 235--248.
\url{https://arxiv.org/abs/2105.06332}.

\bibitem{Gradient}
G. S. H. Cruttwell, B. Gavranovi\'{c}, N. Ghani, P. Wilson, and F. Zanasi.
\emph{Categorical Foundations of Gradient-Based Learning.}
ESOP 2022.
\url{https://arxiv.org/abs/2103.01931}.

\bibitem{SmitheI}
T. St Clere Smithe.
\emph{Compositional Active Inference I: Bayesian Lenses. Statistical Games.}
\url{https://arxiv.org/abs/2109.04461}.

\bibitem{SmitheII}
T. St Clere Smithe.
\emph{Compositional Active Inference II: Polynomial Dynamics. Approximate Inference Doctrines.}
\url{https://arxiv.org/abs/2208.12173}.

\bibitem{Petri}
J. C. Baez and J. Master. \emph{Open Petri Nets.}
Mathematical Structures in Computer Science 30(3) (2020), 314--341.
\url{https://arxiv.org/abs/1808.05415}.

\bibitem{CausalError}
E. F. Rischel and S. Weichwald.
\emph{Compositional Abstraction Error and a Category of Causal Models.}
UAI 2021, Proceedings of Machine Learning Research 161.
\url{https://proceedings.mlr.press/v161/rischel21a.html}.

\bibitem{Contextuality}
S. Abramsky and A. Brandenburger.
\emph{The Sheaf-Theoretic Structure of Non-Locality and Contextuality.}
New Journal of Physics 13 (2011), 113036.
\url{https://arxiv.org/abs/1102.0264}.

\bibitem{HR}
S. Colton, A. Bundy, and T. Walsh.
\emph{Automatic Concept Formation in Pure Mathematics.}
IJCAI 1999.
\url{https://ijcai.org/proceedings/1999/112}.

\bibitem{DreamCoder}
K. Ellis et al.
\emph{DreamCoder: Bootstrapping Inductive Program Synthesis with Wake-Sleep Library Learning.}
PLDI 2021, 835--850.
\url{https://people.csail.mit.edu/asolar/papers/EllisWNSMHCST21.pdf}.

\bibitem{ScientificSchema}
F. Y. Wang and M. J. Buehler.
\emph{Self-Revising Discovery Systems for Science: A Categorical Framework for Agentic Artificial Intelligence.}
arXiv preprint, 2026. Source claims inspected; experiments not independently reproduced.
\url{https://arxiv.org/abs/2606.01444}.

\bibitem{Outgrow}
S. Ma et al.
\emph{Learning to Outgrow a Theory: Experimental Discovery Beyond the Initial Hypothesis Space.}
arXiv preprint, 2026. Source claims inspected; experiments not independently reproduced.
\url{https://arxiv.org/abs/2610.07627}.

\end{thebibliography}
\end{document}



